% ═════════════════════════════════════════════════════════════════════════════
\section{RSA Public-Key Cryptography}
% ═════════════════════════════════════════════════════════════════════════════

RSA is a public-key cryptosystem based on the difficulty of factoring the
product of two large primes. A public key $(n, e)$ is used for encryption;
a private key $d$ is used for decryption. The security of the scheme rests on
the following structure.

% ═════════════════════════════════════════════════════════════════════════════
\section{Background}
% ═════════════════════════════════════════════════════════════════════════════

Choose two distinct primes $p$ and $q$. Set $n = pq$ and $\phi(n) = (p-1)(q-1)$.
Choose $e$ with $\gcd(e, \phi(n)) = 1$ (the \emph{encryption exponent}), and
find $d$ with
\[
    ed \equiv 1 \pmod{\phi(n)}
\]
(the \emph{decryption exponent}). To encrypt a message block $m < n$:
\[
    E(m) = m^e \pmod{n}.
\]
To decrypt a ciphertext block $c$:
\[
    D(c) = c^d \pmod{n}.
\]
By Euler's theorem, $D(E(m)) = m^{ed} \equiv m \pmod{n}$.

% ═════════════════════════════════════════════════════════════════════════════
\section{Finding the Decryption Key}
% ═════════════════════════════════════════════════════════════════════════════

\textbf{Public key:} $n = 11009$, $e = 1507$. It is given that $n = 101 \times 109$.

\subsection{Computing \texorpdfstring{$\phi(n)$}{φ(n)}}

\[
    \phi(n) = (p-1)(q-1) = (101-1)(109-1) = 100 \times 108 = 10800.
\]

\subsection{Verifying \texorpdfstring{$\gcd(e, \phi(n)) = 1$}{gcd(e,φ(n))=1}}

We first confirm that $e = 1507$ and $\phi(n) = 10800$ are coprime, as required.
This is verified by the computation below, which produces remainder $1$.

\subsection{Extended Euclidean Algorithm}

We seek $d$ such that $10800x + 1507d = 1$ (i.e.\ $1507d \equiv 1 \pmod{10800}$).

\textbf{Forward pass (Euclidean algorithm):}
\begin{align*}
    10800 &= 1507 \times 7 + 251, \\
    1507  &= 251  \times 6 + 1, \\
    251   &= 1    \times 251 + 0.
\end{align*}
The last non-zero remainder is $1$, confirming $\gcd(1507, 10800) = 1$.

\textbf{Back-substitution:}
\begin{align*}
    1 &= 1507 - 251 \times 6 \\
      &= 1507 - 6(10800 - 1507 \times 7) \\
      &= 1507 \times 43 + 10800 \times (-6).
\end{align*}
Hence $1507 \times 43 \equiv 1 \pmod{10800}$.

\[
    \boxed{d = 43.}
\]

\textbf{Check:} $1507 \times 43 = 64\,801 = 6 \times 10800 + 1$. \checkmark

% ═════════════════════════════════════════════════════════════════════════════
\section{Decrypting the Message}
% ═════════════════════════════════════════════════════════════════════════════

Bob sends the ciphertext blocks (in 4-digit groups):
\[
    R = 8901 \quad 4422 \quad 5873 \quad 0856 \quad 6130 \quad 7864.
\]

Decryption uses $D(c) = c^{43} \pmod{11009}$ for each block. The resulting
plaintext blocks are converted to ASCII by splitting each 4-digit number into
two 2-digit codes (blank $= 32$, letters $A$--$Z$ map to $65$--$90$).

\begin{center}
\renewcommand{\arraystretch}{1.4}
\begin{tabular}{ccccc}
    \hline
    \textbf{Ciphertext} & \textbf{Decrypted} & \textbf{Split} & \textbf{ASCII} & \textbf{Characters} \\
    \hline
    $8901$ & $6776$ & $67,\;76$ & C, L & CL \\
    $4422$ & $6965$ & $69,\;65$ & E, A & EA \\
    $5873$ & $8232$ & $82,\;32$ & R, \textvisiblespace & R\textvisiblespace \\
    $0856$ & $6583$ & $65,\;83$ & A, S & AS \\
    $6130$ & $3277$ & $32,\;77$ & \textvisiblespace, M & \textvisiblespace M \\
    $7864$ & $8568$ & $85,\;68$ & U, D & UD \\
    \hline
\end{tabular}
\end{center}

\medskip
The decrypted message is:
\[
    \textbf{CLEAR AS MUD.}
\]

\begin{remark}
The modular exponentiations $c^{43} \pmod{11009}$ above were computed using
repeated squaring (e.g.\ $c^{43} = c^{32} \cdot c^8 \cdot c^2 \cdot c^1$).
\end{remark}
